\subsection{A sharper wild lower bound for the character problem}

\begin{proposition}[B-RD28: wild character ambiguity in every higher height]
\label{prop:B-RD28-wild-character}
\status{Proved; independent parent proof and indexing audit passed.}
Let $p=3$, let $e$ be a positive even integer, and put
\[
 s=v_3(e/2),\qquad L=\lfloor\log_3(e/2)\rfloor,\qquad
 N_{\mathrm w}=1+\min\left\{s,
     \left\lfloor\frac{L+1}{2}\right\rfloor,
     \left\lfloor\frac{L+s+1}{4}\right\rfloor\right\}.
\]
For every height $h\geq2$ there are two height-$h$ BT groups
$M,N$ over one complete discretely valued field $K$ of absolute
ramification $e$, with perfect residue field and $\mu_3\subset K$,
such that both groups acquire good reduction,
\[
                 M[3^{N_{\mathrm w}}]\simeq
                 N[3^{N_{\mathrm w}}],
\]
and their manuscript characters differ on a prime-to-three
inertia element. Consequently the least universal B.4 index,
if denoted by $\nu_{\min}(3,e,h)$, satisfies
\[
                       \nu_{\min}(3,e,h)\geq N_{\mathrm w}.
                                                        \tag{B-RD28.1}
\]
In particular the lower bounds at $(e,h)=(24,2)$ and $(216,2)$
are respectively two and three.
\end{proposition}

\begin{proof}
Apply Theorem~\ref{thm:B-RD26-wild} with $n=N_{\mathrm w}$ and
$m=12$. The construction gives an elliptic curve $E/K$ with
full rational $3^n$-torsion and faithful special inertia group
$C_3\rtimes C_4$ on its supersingular good reduction
$Y^2=X^3-X$. By the Weil pairing, $\mu_{3^n}\subset K$.
Set
\[
 M=E[3^\infty]\times(\mathbf Q_3/\mathbf Z_3)^{h-2},
 \qquad N=(\mathbf Q_3/\mathbf Z_3)^h.
\]
These generic BT groups have isomorphic constant $3^n$-torsion.
The second group has good reduction already over $\mathcal O_K$;
the first acquires good reduction over the degree-twelve field
of the construction. Thus both satisfy the manuscript monodromy
condition with a single good-after-extension graded piece.

For the order-four automorphism $\tau$ of the supersingular
elliptic reduction, $\tau^2=[-1]$. Its two eigenvalues on the
rational covariant Dieudonne module are $i,-i$; one may see the
opposite pair either from the degree-one determinant or from
Frobenius interchanging the two primitive fourth roots in
characteristic three. Hence its trace is zero. On $M$ the trace
is $h-2$, whereas on $N$ it is $h$. The period twist in the
corrected Fontaine dictionary has trivial finite inertia action,
so these are exactly the manuscript character values.
The characters therefore differ. Failure at
$3^n=M(n-1)$ gives $\nu_{\min}\geq n$, with precisely the indexing
specified in the manuscript.
\end{proof}

\begin{remark}[Scope of the lower bound]
This improves the tame Honda bound in
Corollary~\ref{cor:B-RD18} in the wild exceptional ranges.
It does not assert that the displayed lower bound is the least
universal index: the latter concerns arbitrary pairs of groups
with isomorphic truncations, and the comparison group need not
be constant. The independent characteristic-two example in
Proposition~\ref{prop:B-RD22} already illustrates that distinction.
The finite upper bound in Theorem~\ref{thm:B-RD17} continues to
apply to every pair satisfying the full filtration condition.
The numerical formula here is included among the $25000$ exact
least-index comparisons in the certificate for
Theorem~\ref{thm:B-RD26-wild}; no new unexecuted computation is
being used.
\end{remark}
